\scanpage{15}
topologie faible) donc \unsure{fblt} \inserted{\unsure{relt}} compact dans $L^1(K)$ ; \deleted{d'où} \inserted{i.e.} \deleted{facilement \unsure{qu'il|que il} est \unsure{fblt} compact dans $L^\infty(K)$.} pour $\sigma(\deleted{\unsure{L^1}}L^\infty(\unsure{K|k})\unsure{,}L^\infty(K))$, donc (étant borné dans $L^\infty(K)$), pour $\sigma(L^\infty(K),L^1(K))$.
% slanted
\margin{\unsure{Comme on|Comme} \unsure{peut|on|va} \unsure{supposer|suppose} $\unsure{a|\omega}$ fini}

Considérons maintenant l'ens $A_{\unsure{P|p}}(K)$ dans $\underline R\times L^\infty(K)$ formé des $(\lambda,f)$ tels que $\lambda\geq P(f)$. $P_2$ signifie que $A_P$ est convexe, $P'_3$ que $A_P$ est \unsure{fblt} fermé, donc (Hahn Banach) $A_P$ est une intersection de demi-espaces \unsure{fblt} fermés,
% slanted
\margin{fait général pour les \unsure{fonct.} convexes semi continues inférieurement} qui sont forcément de la forme $\lambda\geq L_i(f)+\deleted{\unsure{a}}a_i$ ($L_i$ \unsure{formes|forme} \unsure{fblt} \unsure{continues|continue} sur $L^\infty(K)$, i.e. $L_i\in L^1(K)$ ; \deleted{\unsure{$L_i$} \unsure{en vertu de}} $a_i\in\underline R$) d'où \deleted{$P(f)=\unsure{\mathop{\rm Sup}}$} $P(f)=\mathop{\rm Sup}\limits_i(L_i(f)+a_i)$. En vertu de $P_1$ on a \deleted{\unsure{$P(f)$}} aussi, pour \deleted{tout} \deleted{\unsure{$g$}} \inserted{$i$ et \deleted{tout} $g\in L^\infty(K)$} \unsure{spectralt|spectralement} \unsure{équivalent} à $f$, $P(f)\geq L_i(g)+a_i$ d'où \inserted{en vertu de prop. $2$,} faisant varier $g$ $P(f)\geq\int\varphi_{L_i}\varphi_f+a_i$. Posons $\varphi_{L_i}\unsure{:|=}\varphi_i$ ; on a donc
\begingroup\csname tagsleft@true\endcsname
\[
\boxed{P(f)=\mathop{\rm Sup}\limits_{i\in I}a_i+\int_0^{\unsure{\lvert K\rvert}}\varphi_i\varphi_f}
\tag*{(1)}
\]
\endgroup
% archive sheet number: 16

